\documentclass[reqno,11pt]{amsart}
\usepackage{math327-lecture}
\usepackage{needspace}
\usepackage{tikz}
\lectureinfo{19}{Oct. 23, 2026}{Products of Cyclic Groups, Internal Direct Products, and Quotient Groups}
\artinsections{2.11--2.12}



\tcolorboxenvironment{example}{lecturebox,colback=definitiongreen,colframe=definitionborder}
\definecolor{remarkpurple}{HTML}{F2ECFA}
\definecolor{remarkpurpleborder}{HTML}{A68CC7}
\tcolorboxenvironment{remark}{lecturebox,colback=remarkpurple,colframe=remarkpurpleborder}

\begin{document}
\raggedbottom
\makelecturetitle

Last time we introduce the notion of an external direct product of two groups. If $G_1$, and $G_2$ are two groups, this is defined as the set \be{} G_1 \times G_2 \ee equipped with componentwise multiplication. The identity if $(e_{G_1}, e_{G_2})$ and the inverse of $(g_1, g_2)$ is $(g_1^{-1}, g_2^{-1})$.


\section{Products of cyclic groups}

Write $C_d$ for a cyclic group of order $d$.

\begin{lemma}\label{lem:coprime-cyclic}
If $r,s\geq1$ and $\gcd(r,s)=1$, then
\[
  C_{rs}\cong C_r\times C_s.
\]
\end{lemma}

\begin{warningbox}
The coprimality assumption matters. For example,
$C_4\not\cong C_2\times C_2$. Every element $(u,v)$ of
$C_2\times C_2$ satisfies
\[
  (u,v)^2=(u^2,v^2)=(1,1),
\]
whereas $C_4$ has an element of order $4$. An isomorphism preserves
element orders, so these groups cannot be isomorphic.
\end{warningbox}

\begin{proof}[Proof of Lemma~\ref{lem:coprime-cyclic}]
Choose generators
\[
  C_{rs}=\langle x\rangle,\qquad C_r=\langle z\rangle,
  \qquad C_s=\langle y\rangle.
\]
Define
\[
  \varphi:C_{rs}\to C_r,\quad x^i\mapsto z^i,
  \qquad
  \psi:C_{rs}\to C_s,\quad x^i\mapsto y^i.
\]
These maps are well defined: if $x^i=x^j$, then $rs\mid(i-j)$,
so both $r$ and $s$ divide $i-j$, giving $z^i=z^j$ and $y^i=y^j$.
They are homomorphisms, since, for instance,
\[
  \varphi(x^ix^j)=z^{i+j}=z^iz^j.
\]
Define a map using these two coordinates:
\[
  \Phi:C_{rs}\longrightarrow C_r\times C_s,
  \qquad x^i\longmapsto(z^i,y^i).
\]
It is a homomorphism because multiplication is componentwise:
\[
  \Phi(x^ix^j)=(z^{i+j},y^{i+j})
  =(z^i,y^i)(z^j,y^j)=\Phi(x^i)\Phi(x^j).
\]
The two coordinate kernels are
\[
  \ker(\varphi)=\langle x^r\rangle,\qquad
  \ker(\psi)=\langle x^s\rangle.
\]
Consequently, $x^i\in\ker(\Phi)$ if and only if $r\mid i$ and
$s\mid i$. Since $\gcd(r,s)=1$, this is equivalent to $rs\mid i$,
so $\ker(\Phi)=\{1\}$ and $\Phi$ is injective.

Both its domain and its codomain have order $rs$. Thus $\Phi$ is also
surjective, and hence an isomorphism. Equivalently, the counting
formula gives $|\im(\Phi)|=rs$, so its image is all of
$C_r\times C_s$.
\end{proof}

% Cyclic and internal products carried forward from typed Lecture 18.
% Source: handwritten Lecture 13, Oct. 1, 2025, pp. 2--11.
\section{Internal direct products}

We now take the internal
point of view: given subgroups $H,K\leq G$, when can we identify $G$
with $H\times K$ using the multiplication map $(h,k)\mapsto hk$?

\subsection{The product of two subgroups}

Let $H,K\leq G$. The \emph{product set} of $H$ and $K$ is
\[
  HK=\{hk\mid h\in H,\ k\in K\}\subseteq G.
\]
Here $HK$ consists of elements of $G$, whereas $H\times K$ consists
of ordered pairs and has the external direct product group structure.

\begin{warningbox}
The set $HK$ need not be a subgroup of $G$.
\end{warningbox}

\begin{example}
In $S_3$, let $H=\langle(1\,2)\rangle$ and
$K=\langle(2\,3)\rangle$. Composing permutations from right to left,
\[
  HK=\{1,(1\,2),(2\,3),(1\,2\,3)\}.
\]
This set is not closed under multiplication: both $(2\,3)$ and
$(1\,2)$ belong to $HK$, but their product $(1\,3\,2)$ does not.
\end{example}


\begin{lemma}\label{lem:product-subgroup}
If either $H$ or $K$ is normal in $G$, then $HK=KH$ and $HK\leq G$.
\end{lemma}
\begin{proof}
Suppose first that $H\trianglelefteq G$. For $h\in H$ and $k\in K$,
\[
  hk=k(k^{-1}hk)\in KH,\qquad
  kh=(khk^{-1})k\in HK,
\]
because the conjugates $k^{-1}hk$ and $khk^{-1}$ lie in $H$.
Thus $HK=KH$.

The identity lies in $HK$. For closure under multiplication, write
$k_1h_2=h'_2k'_1$ using $KH=HK$. Then
\[
  (h_1k_1)(h_2k_2)=(h_1h'_2)(k'_1k_2)\in HK.
\]
Also, $(hk)^{-1}=k^{-1}h^{-1}\in KH=HK$, so $HK$ is closed under
inverses and hence is a subgroup. If $K\trianglelefteq G$, interchange
$H$ and $K$ in the argument.
\end{proof}

\Needspace{10\baselineskip}
\subsection{Recognizing a direct product}

\begin{define}[Internal direct product]
A group $G$ is the \emph{internal direct product} of subgroups $H$ and
$K$ if
\[
  H\trianglelefteq G,\qquad K\trianglelefteq G,\qquad
  H\cap K=\{1\},\qquad HK=G.
\]
\end{define}

Let's begin with a simple lemma.


\begin{lemma}[Unique factorization]\label{lem:unique-factorization}
	If $H\cap K=\{1\}$, the multiplication map of sets
	\[
	p:H\times K\longrightarrow G,\qquad (h,k)\longmapsto hk
	\]
	is injective. Thus each element of $HK$ has a unique expression $hk$
	with $h\in H$ and $k\in K$.
\end{lemma}
\begin{proof}
	If $h_1k_1=h_2k_2$, then
	\[
	h_2^{-1}h_1=k_2k_1^{-1}.
	\]
	The left side lies in $H$ and the right side lies in $K$. Their common
	value therefore belongs to $H\cap K=\{1\}$, so $h_1=h_2$ and $k_1=k_2$.
\end{proof}


\begin{theorem}[Recognizing direct products]\label{thm:internal-product}
If $G$ is the internal direct product of $H$ and $K$, then
\[
  G\cong H\times K.
\]
The isomorphism sends $g=hk$ to $(h,k)$.
\end{theorem}
\begin{proof}
Since $HK=G$, every $g\in G$ can be written $g=hk$. By
Lemma~\ref{lem:unique-factorization}, this expression is unique. Hence
\[
  \pi:G\longrightarrow H\times K,\qquad hk\longmapsto(h,k)
\]
is well defined.

\emph{First we show that elements of $H$ commute with elements of $K$.}
For $h\in H$ and $k\in K$, consider the commutator
\[
  [h,k]=hkh^{-1}k^{-1}.
\]
Normality of $K$ gives $hkh^{-1}\in K$, so $[h,k]\in K$.
Normality of $H$ gives $kh^{-1}k^{-1}\in H$, so
$[h,k]=h(kh^{-1}k^{-1})\in H$. Consequently,
\[
  [h,k]\in H\cap K=\{1\},\qquad\text{and thus }hk=kh.
\]

Now let $g_1=h_1k_1$ and $g_2=h_2k_2$. Since $k_1h_2=h_2k_1$,
\[
  g_1g_2=h_1k_1h_2k_2=(h_1h_2)(k_1k_2).
\]
It follows that
\[
  \pi(g_1g_2)=(h_1h_2,k_1k_2)
  =(h_1,k_1)(h_2,k_2)=\pi(g_1)\pi(g_2).
\]
Thus $\pi$ is a homomorphism. It is bijective, with inverse
$p(h,k)=hk$, and hence is an isomorphism.
\end{proof}

\begin{exercise}
Verify directly that the maps $\pi$ and $p$ in
Theorem~\ref{thm:internal-product} are inverses. Then check that the
subgroups $\widetilde G_1$ and $\widetilde G_2$ of the external product
$G_1\times G_2$ satisfy all four conditions in the definition of an
internal direct product.
\end{exercise}

\begin{remark}
Both normality assumptions matter. If only $H\trianglelefteq G$ is
assumed, together with $H\cap K=\{1\}$ and $HK=G$, the resulting
structure is called a \emph{semidirect product}, written
$G\cong H\rtimes K$; the factors need not commute.

For example, in $S_3$ take $H=\langle a\rangle$ with $a=(1\,2\,3)$
and $K=\langle b\rangle$ with $b=(1\,2)$. Then $H\trianglelefteq S_3$,
$H\cap K=\{1\}$, and
\[
  HK=\{1,a,a^2,b,ab,a^2b\}=S_3.
\]
But $K$ is not normal: $aba^{-1}=(2\,3)\notin K$.
Moreover, $bab^{-1}=a^{-1}\ne a$, so the factors do not commute.
Thus $S_3\cong C_3\rtimes C_2$ is a semidirect product; it is not
isomorphic to the abelian group $C_3\times C_2$.
\end{remark}


\section{Quotient groups}

Let $H\leq G$. Recall that $G/H$ denotes the \emph{set} of left
cosets $gH$. We would like to multiply cosets by the rule
\[
  (g_1H)(g_2H)=(g_1g_2)H.
\]
Before checking the group axioms, we must check that this rule does not
depend on the representatives $g_1$ and $g_2$.

\subsection{Why normality is needed}

If $g'_1=g_1h_1$ and $g'_2=g_2h_2$ with $h_1,h_2\in H$, then
$g'_1H=g_1H$ and $g'_2H=g_2H$. But
\[
  g'_1g'_2=g_1h_1g_2h_2
  =g_1g_2\bigl(g_2^{-1}h_1g_2\bigr)h_2.
\]
If $H\trianglelefteq G$, the parenthesized element lies in $H$, and
therefore $g'_1g'_2H=g_1g_2H$, as required.

\begin{example}[Failure for a nonnormal subgroup]
Take $G=S_3$ and $H=\{1,(1\,2)\}$. The identity coset has both
$1$ and $(1\,2)$ as representatives. Using $g_2=(2\,3)$ gives
\[
  1(2\,3)H=(2\,3)H,\qquad
  (1\,2)(2\,3)H=(1\,2\,3)H.
\]
These cosets are distinct: otherwise
$(2\,3)^{-1}(1\,2\,3)=(1\,3)$ would belong to $H$.
Thus the proposed product depends on the chosen representative.
\end{example}

\begin{remark}
In fact, normality is necessary for this multiplication rule to be well
defined. If it is well defined, then for every $h\in H$ and $g\in G$,
\[
  (hH)(gH)=(1H)(gH),\qquad\text{so }hgH=gH.
\]
Hence $g^{-1}hg\in H$. Applying this for every $g$ proves
$H\trianglelefteq G$.
\end{remark}

\subsection{The quotient construction}

\begin{theorem}[Quotient group]
Let $N\trianglelefteq G$. The set $G/N$ is a group under
\[
  (gN)(hN)=(gh)N.
\]
Its identity is $N=1N$, and the inverse of $gN$ is $g^{-1}N$.
\end{theorem}
\begin{proof}
The preceding calculation proves that multiplication is well defined.
Associativity follows from associativity in $G$:
\[
 ((gN)(hN))(kN)=((gh)k)N=(g(hk))N=(gN)((hN)(kN)).
\]
Also, $N(gN)=(gN)N=gN$, and
\[
  (gN)(g^{-1}N)=(g^{-1}N)(gN)=N.
\]
These are the group axioms.
\end{proof}

\Needspace{10\baselineskip}
\begin{theorem}[The quotient map]
The map
\[
  \pi:G\longrightarrow G/N,\qquad g\longmapsto gN
\]
is a surjective homomorphism with kernel $N$.
\end{theorem}
\begin{proof}
By definition of the quotient multiplication,
$\pi(gh)=(gh)N=(gN)(hN)=\pi(g)\pi(h)$.
Every coset has a representative, so $\pi$ is surjective. Finally,
\[
  g\in\ker\pi\quad\Longleftrightarrow\quad gN=N
  \quad\Longleftrightarrow\quad g\in N.
\]
\end{proof}

\begin{corollary}
If $G$ is finite, then $|G/N|=[G:N]=|G|/|N|$.
\end{corollary}
\begin{proof}
The elements of $G/N$ are precisely the left cosets of $N$ in $G$.
Each has $|N|$ elements, and they partition $G$.
\end{proof}

\begin{example}
In additive notation, $n\Z\trianglelefteq\Z$ and
\[
  (a+n\Z)+(b+n\Z)=(a+b)+n\Z.
\]
Thus $\Z/n\Z$ is the familiar group of integers modulo $n$.
The quotient map sends an integer to its residue class; its kernel is
$n\Z$.
\end{example}

\begin{exercise}
For the external product $H\times K$, show that
\[
  (H\times K)/(H\times\{1\})\cong K
\]
by defining a map $(h,k)(H\times\{1\})\mapsto k$.
Check that it is well defined and is a bijective homomorphism.
\end{exercise}

\end{document}
